\par\Needspace{6\baselineskip}\medskip
\hypertarget{native-143-article104}{}\pdfbookmark[2]{104.}{article-104}
\begin{center}104.\end{center}
It remains to determine the elements themselves from $z$, $n$, or $c$. Setting
\[
\alpha\BareRoot{(ee-1)}=\beta,
\]
\noindent we obtain from equation 6 of Art. 99
\[
\mathrm{18)}\qquad\beta=\frac{\sin f\mathpunct{.}\BareRoot{rr'}}{\operatorname{tang}2n}\mathpunct{\cdot}
\]
\noindent Combining this formula with 12 and $12^*$ of Art. 99 gives
\[
\begin{aligned}
\mathrm{19)}\qquad&\BareRoot{(ee-1)}=\operatorname{tang}\psi=\frac{\operatorname{tang}f\operatorname{tang}2n}{2(l-z)}\\[.6em]
\mathrm{19^*)}\qquad&\operatorname{tang}\psi=-\frac{\operatorname{tang}f\operatorname{tang}2n}{2(L+z)},
\end{aligned}
\]
\noindent From these formulas the eccentricity can be calculated conveniently and accurately. Dividing $\beta$ by $\BareRoot{(ee-1)}$ gives $\alpha$, and multiplying them gives $p$. Thus we have
\[
\begin{aligned}
\alpha&=\frac{2(l-z)\cos f\mathpunct{.}\BareRoot{rr'}}{\operatorname{tang}2n^2}
=\frac{2mm\cos f\mathpunct{.}\BareRoot{rr'}}{yy\operatorname{tang}2n^2}\\[.7em]
&=\frac{kktt}{4yyrr'\cos f^2\operatorname{tang}2n^2}\\[.7em]
&=\frac{-2(L+z)\cos f\mathpunct{.}\BareRoot{rr'}}{\operatorname{tang}2n^2}
=\frac{-2MM\cos f\mathpunct{.}\BareRoot{rr'}}{YY\operatorname{tang}2n^2}\\[.7em]
&=\frac{kktt}{4YYrr'\cos f^2\operatorname{tang}2n^2}
\end{aligned}
\]
\[
\begin{aligned}
p&=\frac{\sin f\mathpunct{.}\operatorname{tang}f\mathpunct{.}\BareRoot{rr'}}{2(l-z)}
=\frac{yy\sin f\mathpunct{.}\operatorname{tang}f\mathpunct{.}\BareRoot{rr'}}{2mm}\\[.7em]
&=\left(\frac{yrr'\sin2f}{kt}\right)^2\\[.7em]
&=\frac{-\sin f\mathpunct{.}\operatorname{tang}f\mathpunct{.}\BareRoot{rr'}}{2(L+z)}
=\frac{-YY\sin f\mathpunct{.}\operatorname{tang}f\mathpunct{.}\BareRoot{rr'}}{2MM}\\[.7em]
&=\left(\frac{Yrr'\sin2f}{kt}\right)^2\mathpunct{\cdot}
\end{aligned}
\]
\noindent\SourceLeaf{144}The third and sixth expressions for $p$, which are identical to formulas 18 and $18^*$ of Art. 95, show that what was stated there about the meaning of the quantities $y$, $Y$ also applies to the hyperbola.
\par
Combining equations 6 and 9 of Art. 99 yields
\[
(r'-r)\BareRoot{\frac{ee-1}{rr'}}=e\sin f\mathpunct{.}\left(C-\frac1C\right);
\]
\noindent so that, introducing $\psi$ and $\omega$ and setting $C=\operatorname{tang}(45^0+N)$, we get
\[
\mathrm{20)}\qquad\operatorname{tang}2N=\frac{2\sin\psi\operatorname{tang}2\omega}{\sin f\cos2\omega}\mathpunct{\cdot}
\]
\noindent Once $C$ has been found in this way, we obtain for both positions the values of the quantity denoted by $u$ in Art. 21. Equation III of Art. 21 then gives
\[
\begin{aligned}
\operatorname{tang}\tfrac12v&=\frac{C-c}{(C+c)\operatorname{tang}\tfrac12\psi}\\[.6em]
\operatorname{tang}\tfrac12v'&=\frac{Cc-1}{(Cc+1)\operatorname{tang}\tfrac12\psi}
\end{aligned}
\]
\noindent or, replacing $C$, $c$ with the angles $N$, $n$,
\[
\begin{aligned}
\mathrm{21)}\qquad&\operatorname{tang}\tfrac12v=\frac{\sin(N-n)}{\cos(N+n)\operatorname{tang}\tfrac12\psi}\\[.6em]
\mathrm{22)}\qquad&\operatorname{tang}\tfrac12v'=\frac{\sin(N+n)}{\cos(N-n)\operatorname{tang}\tfrac12\psi}\mathpunct{\cdot}
\end{aligned}
\]
\noindent From this we determine the true anomalies $v$, $v'$. Comparing their difference with $2f$ also serves to check the calculation.
\par
Finally, formula XI of Art. 22 readily gives the time interval from perihelion to the time corresponding to the first position as
\[
=\frac{\alpha^{\frac32}}k\left\{\frac{2e\cos(N+n)\sin(N-n)}{\cos2N\cos2n}-\log\operatorname{hyp}\frac{\operatorname{tang}(45^0+N)}{\operatorname{tang}(45^0+n)}\right\}
\]
\noindent and similarly gives the time interval from perihelion to the time corresponding to the second position as
\[
=\frac{\alpha^{\frac32}}k\left\{\begin{aligned}&\frac{2e\cos(N-n)\sin(N+n)}{\cos2N\cos2n}\\[.5em]&-\log\operatorname{hyp}\operatorname{tang}(45^0+N)\operatorname{tang}(45^0+n)\end{aligned}\right\}.
\]
\noindent If, then, the first time is set $=T-\tfrac12t$ and the second is therefore $=T+\tfrac12t$, we obtain
\[
\mathrm{23)}\qquad T=\frac{\alpha^{\frac32}}k\left\{\frac{e\operatorname{tang}2N}{\cos2n}-\log\operatorname{hyp}\operatorname{tang}(45^0+N)\right\},
\]
\noindent\SourceLeaf{145}which gives the time of passage through perihelion. Finally,
\[
\mathrm{24)}\qquad t=\frac{2\alpha^{\frac32}}k\left\{\frac{e\operatorname{tang}2n}{\cos2N}-\log\operatorname{hyp}\operatorname{tang}(45^0+n)\right\},
\]
\noindent This equation may, if desired, be used as a final check on the calculation.
